\documentclass[11pt,a4paper]{article}
% Shared preamble for RuPhO English problem statements (match T1 2026 style).
% Use: \input{latex/rupho_en_preamble.tex} after \documentclass.
\usepackage[margin=1in]{geometry}
\usepackage{amsmath,amssymb}
\usepackage{graphicx}
\usepackage{float}
\usepackage{enumitem}
\usepackage{microtype}
\usepackage{caption}
\usepackage{tabularx}
\usepackage{hyperref}

\captionsetup{font=small,labelfont=bf,skip=6pt}

\setlist[itemize]{leftmargin=1.2cm}
\setlist[enumerate]{leftmargin=1.2cm}

% One outer box per subproblem: [id | statement | points] (no inner rules)
\newcommand{\problembox}[3]{%
  \par\addvspace{0.42em}%
  \noindent
  \setlength{\fboxsep}{7.5pt}%
  \setlength{\fboxrule}{0.95pt}%
  \fbox{%
    \begin{minipage}{\dimexpr\linewidth-2\fboxsep-2\fboxrule\relax}
    \begin{tabularx}{\linewidth}{@{}>{\raggedright\arraybackslash\bfseries}p{2.1em} @{\hspace{0.9em}} >{\raggedright\arraybackslash}X @{\hspace{0.9em}} >{\raggedleft\arraybackslash\bfseries}p{2.4em}@{}}
    #1 & #2 & #3
    \end{tabularx}
    \end{minipage}%
  }%
  \par\addvspace{0.32em}%
}
\title{\textbf{Sun spots --- Solution}}
\author{\normalsize X22 (2022) --- English translation}
\date{}
\begin{document}
\maketitle
\problembox{A1}{You are presented with a table of experimental data $\omega(\varphi)$. The data was obtained based on the Doppler effect, which affects the frequency of electromagnetic radiation from the Sun recorded on Earth. Define parameters $a$ and $b$.\\ Let's build a graph $\omega(\sin^2 \varphi)$. We draw an approximating line along the initial section, from which we find the parameters $a$ and $b$.\\ Answer: $a = (2.99 \pm 0.2) \cdot 10^{-6}~1/\text{s}$}{2.00}
\begin{figure}[H]
    \centering
    \includegraphics[width=0.72\textwidth]{PE_S_EN_assets/fig1.jpg}
\end{figure}
\problembox{A2}{Find the limits of applicability of the formula given above.\\ From the graph we find the angle $\varphi_{\text{max}}$ after which the points significantly deviate from the straight line.\\ Answer: $\varphi_{\text{max}} = 50^{\circ}$}{0.30}
\problembox{B1}{Obtain all components of the magnetic field on the surface of the Sun $B_{0,r}(\phi)$, $B_{0,\phi}(\phi)$, $B_{0,\lambda} (\phi)$. Express your answers using $\varphi$, $\mathfrak{M}$, $R_0$ and $\mu_0$.\\ Let us write the field of the magnetic dipole at latitude $\varphi$: \[ \vec{B}(\varphi) = \frac{\mu_0}{4\pi} \frac{\mathfrak{M}}{R_0^3} (3\vec{n} \sin \varphi - \vec{z} ) \] Projecting this onto the axes corresponding to $r$, $\varphi$ and $\lambda$ we obtain:\\ Answer: \[B_{0,r} (\varphi) = \dfrac{\mu_0}{4\pi} \dfrac{\mathfrak{M}}{R_0^3} 2 \sin \varphi\]}{0.30}
\problembox{B2}{Obtain the magnitude of the magnetic field $B_e$ at the equator inside the Sun near its surface. Express your answer using $\mathfrak{M}$, $\mu_0$, $R_0$, $\xi$.\\ Let's choose a surface stretched over magnetic lines. Then the magnetic field flux through the equatorial section inside the Sun is equal to the flux through one of the hemispheres.\\ Thus $2\pi R_0 H_0 B_e = \Phi_0$ means $B_e = \dfrac{\Phi_0}{2\pi \xi R_0^2}=53~\text{T}$. On the other hand, we can calculate the flux $\Phi_0$ by integrating the dipole field: \[ \Phi_0 = \int\limits_0^{\pi/2} B_{0,r} 2 \pi R_0^2 \cos \varphi d\varphi = \frac{\mu_0 \mathfrak{M}}{R_0} \int\limits_0^{\pi/2} \sin \varphi \cos \varphi d\varphi = \frac{\mu_0 \mathfrak{M}}{2R_0}. \] Substitute this into the previous formula and get the following:\\ Answer: \[ B_e = \frac{\mu_0 \mathfrak{M}}{4\pi\xi R_0^3}. \]}{0.60}
\begin{figure}[H]
    \centering
    \includegraphics[width=0.72\textwidth]{PE_S_EN_assets/fig2.jpg}
\end{figure}
\problembox{B3}{Find the numerical value of $B_e$.\\ Answer: \[B_e = \dfrac{\Phi_0}{2\pi \xi R_0^2}=53~\text{T}\]}{0.10}
\problembox{B4}{Obtain the magnitude of the magnetic field $B_\varphi(\varphi)$ inside the Sun at its surface as a function of latitude $\varphi$. Express your answer using $\varphi$, $B_e$.\\ Now we will also select a surface stretched over magnetic lines. Only now it will extend from the equator to the desired latitude $\varphi$.\\ Then to calculate $B_\varphi$ you need to find the flux $\Phi(\varphi)$ of the magnetic field of the dipole through the part of latitude from $0$ to $\varphi$. This is the integral already familiar to us: \[\Phi(\varphi) (\varphi)= \int\limits_0^{\varphi} B_{0,r} 2 \pi R_0^2 \cos \varphi d\varphi = \frac{\mu_0 \mathfrak{M}}{2R_0} \sin^2 \varphi. \] The flux through the surface is equal to 0: \[ \Phi(\varphi) + 2\pi R_0 \cos \varphi \, H_0 \cos^2 \varphi\, B_\varphi(\varphi) = \Phi_0, \] then: \[ 2\pi\xi R_0^2 \cos^3 \varphi \,B_\varphi (\varphi) = \Phi_0 \cos^2 \varphi. \] So\\ Answer: \[ B_\varphi = \frac{B_e}{\cos \varphi}.\]}{0.60}
\begin{figure}[H]
    \centering
    \includegraphics[width=0.72\textwidth]{PE_S_EN_assets/fig3.jpg}
\end{figure}
\problembox{C1}{Precisely calculate 30 pairs of points $(\varphi, \lambda)$ through which the magnetic line piercing points $\varphi=\pm 50^\circ$, $\lambda=0$ will pass at time $t=150~\text{}$.\\ Let the magnetic line pass through the points $(\pm\varphi_0,\lambda_0)$, where $\varphi_0=50^\circ$, $\lambda_0=0^\circ$. Then its equation in angular coordinates is given as follows: \[ \lambda = \omega(\varphi)t = (a - b\sin^2 \varphi)t \]}{1.70}
\problembox{C2}{Using this data, plot on your answer sheet what this line will look like for an earthly observer.\\ Points with angular coordinates in the plane have the following coordinates: \[ \begin{cases} x = R \cdot \sin \lambda\, \cos \varphi\\ y = R \cdot \sin \varphi \end{cases} \] Let’s draw a line through the recalculated points:}{1.00}
\begin{figure}[H]
    \centering
    \includegraphics[width=0.72\textwidth]{PE_S_EN_assets/fig4.jpg}
\end{figure}
\problembox{C3}{Find the angle $\psi (\varphi,t)$ that the magnetic lines make with the meridian. Express your answer using $\varphi$, $t$, $a$, $b$.\\ From the picture it is clear that \[ \tan \psi = \frac{d \lambda}{d \varphi} \cos \varphi = 2bt \sin \varphi \cos^2 \varphi, \] is from\\ Answer: \[ \psi = \arctan \left( 2bt \sin \varphi \, \cos^2 \varphi \right). \]}{0.30}
\begin{figure}[H]
    \centering
    \includegraphics[width=0.72\textwidth]{PE_S_EN_assets/fig5.jpg}
\end{figure}
\problembox{4}{Let the magnetic lines inside the Sun near its surface make an angle $\psi$ with the meridian. Find the magnitude of the magnetic field $B$ at this point. Express your answer using $\psi$, $\varphi$, $B_e$.\\ The magnetic field flux through the surfaces discussed in Part B should not change, so $B_\varphi=\text{const}$. Then \[ B = \frac{B_\varphi}{\cos \psi} \] and therefore\\ Answer: \[\frac{B_e}{\cos \varphi \, \cos \psi} \]}{0.20}
\problembox{D1}{Using the $\sin \psi \approx 1 $ approximation, find $B_\text{}$.\\ In the approximation $\sin \psi \approx 1$ it is true that $\tan \psi = 1/\cos \psi$, therefore: \[ B = \frac{B_e \tan \psi}{\cos \varphi} = 2B_e bt \sin \varphi, \] means the type of curve, $\varphi(t)$, on which spots appear on the Sun is described by the following expression: \[ B_\text{} = 2B_e bt \sin \varphi, \] i.e. its linearization is as follows: \[ t = \frac{A}{\sin 2\varphi} , \] where $A = \frac{B_{\text{}}}{B_eb}$ and $t$ are counted from the beginning of the 11-year solar cycle (one “butterfly” on the Maunder diagram). By constructing the envelope on the Maunder diagram, we obtain the dependence $t(\sin 2\varphi)$ for critical points (appearance of spots). From it we find $A$, and then $B_{}$.\\ Answer: \[ B_{} = (15 \pm 5) \, T \]}{2.70}
\begin{figure}[H]
    \centering
    \includegraphics[width=0.72\textwidth]{PE_S_EN_assets/fig6.jpg}
\end{figure}
\problembox{D2}{Determine the period of solar activity $T_0$, i.e. the time after which the Sun's magnetic field returns to its original state.\\ We find the period as the time between two points with the same “phase” divided by the number of cycles: $T = 11~\text{}$. Next, this number must be increased by 2 times, since after one cycle the field changes direction.\\ Answer: \[ T_0 = 22 \, year \]}{0.20}
\vspace{1em}
\noindent\textit{Source:} \url{https://pho.rs/p/506/s}
\end{document}